\section{Experimental Setup}
\label{sec:design}

Validation proceeds in three stages, addressing in turn whether the
method works, whether it works on instances it was not designed against,
and whether what works is the adaptivity or merely the extra budget.

\paragraph{Stage 1 --- ablation on the canonical grid.} We reuse the
144-cell combined-stress grid of Sec.~\ref{sec:combined} (72 parameter
combinations $\times$ 2 repeats, $N=21$, $a=19$, $r=6$) and run all five
arms on every cell. Each trial's seed is derived deterministically from
its grid coordinates, so a given cell produces identical sampled counts
under every arm and only the classical stage differs --- the pairing that
makes this a within-subjects design.

\paragraph{Stage 2 --- pre-specified held-out instances.} The grid above
is the one that identified budget starvation as the failure mode, so it
cannot speak to generalisation. We re-ran the identical design, unchanged,
on three instances chosen by a rule fixed before any trial of this stage:
the lowest valid base achieving each target order, excluding the
canonical $a=19$. This gives $(N{=}15,a{=}2,r{=}4)$ --- new modulus and
order; $(N{=}21,a{=}8,r{=}2)$ --- canonical modulus, new order; and
$(N{=}21,a{=}2,r{=}6)$ --- canonical order, new base.

The choice set is small, and we state why rather than presenting it as
preference. Within the ${\le}5$ work-qubit ceiling forced by the
compilation cliff (Sec.~\ref{sec:limitations}, Limitation~1), only
$N\in\{15,21\}$ admits an odd two-prime modulus with an even, non-trivial
order, which factor extraction requires. Since $r=6$ at $N=21$ is already
the largest reachable even order, held-out instances broaden \emph{which}
instance is tested but cannot supply one that is \emph{harder}. They
close the circularity objection --- the method was never tuned on them ---
while only narrowing, not eliminating, the external-validity question.
The selection rule was fixed before any trial of this stage ran and is
documented in the module docstring of the script that ran it. We describe
the instances as \emph{pre-specified}: the rule preceded the data, but
it carries no third-party timestamp, so the reader is being asked to
accept a repository artifact rather than an independent record. We make
no claim to the procedural guarantees of a formal registration.

Before Stage 2 ran, the canonical grid was re-executed with resource
instrumentation attached, gated on reproducing the already-recorded
per-trial outcomes bit-for-bit. It did, on all 144 trials.

\paragraph{Stage 3 --- matched-resource controls.} Each module succeeds
partly because it may spend more than baseline. To separate the
allocation \emph{rule} from the allocation \emph{size}, the instrumented
run records each arm's realised budget per trial --- per-window candidate
counts $m_w$, realised beam width $P'$, final per-window shots $S_w$. For
every trial we then build a control given that same realised budget
\emph{uniformly}, evaluated through the Baseline code path with all three
flags disabled, so no confidence guidance appears anywhere in the
control. The shots-matched control redistributes the total extra shots
evenly across windows; the candidates-matched control uses a fixed
$k=\mathrm{round}(\overline{m_w})$; the beam-matched control uses that
trial's realised $P'$ unconditionally. Only the beam-matched control is
exact; the other two approximate a non-uniform allocation by its mean,
and \SItables{} quantifies the residual and its direction.

These controls are deliberately favourable to the baseline: they receive
a budget that could only be known \emph{after} watching the adaptive rule
choose it. They are oracles, not competitors, and the comparison asks
whether confidence guidance beats an allocator that already knows the
answer it would have given.

\paragraph{All three stages are noiseless.} Every trial reported here was
sampled from an ideal simulator with no noise model. This is deliberate:
budget starvation is the failure mode under test, and injecting noise
would add a second, independently varying source of failure the modules
were not built to address. The cost is that the 27\% of
characterisation-campaign failures attributed to noise
(Fig.~\ref{fig:pareto}) fall outside the evaluated regime, and we make no
claim about behaviour when noise rather than budget is the binding
constraint (Limitation~1).

\paragraph{Outcome.} The primary outcome is binary end-to-end success:
the returned phase estimate yields the true order and a non-trivial
factorisation of $N$. Secondary outcomes (realised candidate count, beam
width, additional shots) are in \SItables{}.
